\documentclass[a4paper,12pt]{article}
\usepackage[utf8]{inputenc}
\usepackage[T1]{fontenc}
\usepackage[polish]{babel}
\usepackage{geometry}
\usepackage{array}
\usepackage{booktabs}
\usepackage{tabularx}
\usepackage{hyperref}
\usepackage{float}

\usepackage{tikz}
\usetikzlibrary{shapes.geometric, positioning, arrows.meta, fit, backgrounds}

\geometry{margin=2.5cm}

\hypersetup{
  colorlinks=true,
  linkcolor=black,
  urlcolor=black,
  citecolor=black
}

% Polecenie do oznaczania miejsca na diagram
\newcommand{\diagrambox}[2][3cm]{%
  \vspace{0.5cm}
  \noindent\fbox{%
    \begin{minipage}{\dimexpr\linewidth-2\fboxsep-2\fboxrule\relax}
    \centering
    \vspace{#1}
    \textit{[Miejsce na diagram: #2]}
    \vspace{#1}
    \end{minipage}%
  }\par
  \vspace{0.5cm}
}

\title{%
  \textbf{System automatycznego paczkomatu}\\[0.5em]
  \large Dokumentacja techniczna%
}
\author{Karol Wziątek, Kamil Mendyk}
\date{Marzec 2026}

\begin{document}

\maketitle
\newpage

\tableofcontents
\newpage

%%%%%%%%%%%%%%%%%%%%%%%%%%%%%%%%%%%%%%%%%%%%%%%%%%
\section{Opis ogólny systemu}

\subsection{Nazwa systemu}
System automatycznego paczkomatu do nadawania i odbioru przesyłek.

\subsection{Opis ogólny}
System paczkomatu umożliwia użytkownikom nadawanie i odbieranie przesyłek
za pomocą automatycznych skrytek dostępnych całodobowo. System współpracuje
z kurierem oraz serwisantem i obsługuje dodatkowe funkcjonalności, takie jak
płatność przy odbiorze.

\subsection{Funkcjonalności}

\subsubsection{Funkcjonalności podstawowe}
\begin{itemize}
  \item odbiór paczki przez użytkownika (wprowadzenie kodu lub skan QR, otwarcie skrytki, potwierdzenie odbioru)
  \item nadanie paczki (wybór skrytki, generowanie kodu nadania, umieszczenie paczki)
  \item dostarczenie paczki przez kuriera (identyfikacja, otwieranie skrytek, umieszczanie paczek)
\end{itemize}

\subsubsection{Funkcjonalności wtórne}
\begin{itemize}
  \item płatność przy odbiorze (obsługa terminala płatniczego, blokada wydania paczki bez płatności)
  \item wysyłanie powiadomień do użytkownika
  \item przechowywanie paczek przez określony czas
  \item obsługa błędnych prób odbioru
  \item obsługa braku dostępnych skrytek
  \item autoryzacja użytkownika i kuriera
  \item możliwość otwarcia wielu skrytek na raz (dla kuriera)
  \item raportowanie dostaw
\end{itemize}

\subsubsection{Funkcjonalności serwisowe}
\begin{itemize}
  \item tryb serwisowy
  \item restart systemu
  \item monitorowanie stanu urządzenia
\end{itemize}

%%%%%%%%%%%%%%%%%%%%%%%%%%%%%%%%%%%%%%%%%%%%%%%%%%
\section{Wymagania systemu}

\subsection{Co system robi?}
System umożliwia odbiór i nadawanie przesyłek, przechowuje je w skrytkach,
obsługuje płatności przy odbiorze oraz komunikuje się z systemem centralnym.
Zarządza dostępem do skrytek i rejestruje operacje.

\subsection{Jak system działa?}

\noindent\textbf{Wejścia:}
\begin{itemize}
  \item kod odbioru lub QR
  \item dane z aplikacji mobilnej
  \item dane wprowadzone przez kuriera
  \item sygnały z czujników skrytek
  \item dane z terminala płatniczego
\end{itemize}

\noindent\textbf{Wyjścia:}
\begin{itemize}
  \item otwieranie i zamykanie skrytek
  \item komunikaty na ekranie
  \item potwierdzenia operacji
  \item komunikacja z systemem centralnym
\end{itemize}

\noindent\textbf{Zależności:}
\begin{itemize}
  \item skrytka otwiera się po poprawnym kodzie lub po dokonaniu płatności
  \item brak wolnych skrytek uniemożliwia nadanie paczki
  \item błędny kod blokuje dostęp
\end{itemize}

\subsection{Środowisko działania}
System jest wbudowany w publicznie dostępne paczkomaty, które mogą być
stawiane zarówno w budynku, jak i na zewnątrz.

\noindent Ryzyka, które należy wziąć pod uwagę:
\begin{itemize}
  \item warunki atmosferyczne np.\ deszcz, śnieg
  \item próba uszkodzenia lub włamanie
  \item brak lub słabe połączenie z serwerem -- skutkuje ograniczeniem lub uniemożliwieniem korzystania z systemu
  \item ataki na serwer
\end{itemize}

\noindent Wymagania techniczne (wynikające głównie z poprzedniego punktu):
\begin{itemize}
  \item wodoodporność urządzenia
  \item bezpieczeństwo systemu zarówno od strony fizycznej, jak i software'u
  \item zabezpieczenie połączenia między serwerem oraz paczkomatem, aby zminimalizować podatność na uszkodzenia
  \item zabezpieczenie serwera przed infiltracją
\end{itemize}

\subsection{Z kim system wchodzi w interakcję?}

\noindent\textbf{Użytkownik:}
\begin{itemize}
  \item odbiera i nadaje paczki
  \item może dokonywać płatności
\end{itemize}

\noindent\textbf{Kurier:}
\begin{itemize}
  \item dostarcza paczki do paczkomatu adresata
  \item odbiera nadane paczki
\end{itemize}

\noindent\textbf{Serwisant:}
\begin{itemize}
  \item naprawia i diagnozuje system
  \item ma dostęp do trybu technicznego
\end{itemize}

\noindent\textbf{System centralny:} jego zadaniem jest m.in.\ monitorowanie stanu paczko\-matów i generowanie kodów odbioru.
\begin{itemize}
  \item paczkomat odbiera od systemu centralnego informacje o paczkach, które znajdują się w paczkomacie
  \item paczkomat wysyła do systemu centralnego informacje o odbiorach i nadaniach paczek
\end{itemize}

\subsection{Wymagania prawne}
\begin{itemize}
  \item ochrona danych osobowych
  \item bezpieczeństwo płatności
  \item normy bezpieczeństwa urządzeń
\end{itemize}

%%%%%%%%%%%%%%%%%%%%%%%%%%%%%%%%%%%%%%%%%%%%%%%%%%
\begin{center}
\resizebox{\textwidth}{!}{%
\begin{tikzpicture}[
    transform shape, 
    scale=0.8, % Adjust this value (e.g., 0.7 or 0.9) to fit
    >=Latex,
    usecase/.style={ellipse, draw, thick, minimum width=3.8cm, minimum height=1.4cm, align=center, fill=blue!5, text width=3.4cm},
    actor/.style={rectangle, align=center, font=\bfseries},
    include/.style={dashed, -{Latex}, thick},
    extend/.style={dashed, -{Latex}, thick},
    assoc/.style={-, thick}
]

% ==========================================
% AKTORZY
% ==========================================
\node[actor] (User) at (-3, 0) {Użytkownik};
\node[actor] (Kurier) at (-3, -9.5) {Kurier};
\node[actor] (Serwisant) at (-3, -15) {Serwisant};

\node[actor] (SysPlat) at (19, 4) {System\\Płatniczy};
\node[actor] (SysCentr) at (19, -5) {System\\Centralny};

% ==========================================
% SYSTEM PACZKOMATU (Ramka)
% ==========================================
\begin{scope}[local bounding box=system]
    
    % --- GŁÓWNE PRZYPADKI UŻYCIA (Oś X = 8) ---
    \node[usecase] (PU4) at (8, 4) {PU4: Płatność\\przy odbiorze};
    \node[usecase] (PU1) at (8, 0) {PU1: Odbiór paczki};
    \node[usecase] (PU5) at (8, -2.5) {PU5: Wysyłanie\\powiadomień};
    \node[usecase] (PU2) at (8, -5) {PU2: Nadanie paczki};
    \node[usecase] (PU3) at (8, -8) {PU3: Dostarczenie\\przez kuriera};
    
    % Nowy Przypadek Użycia (PU8)
    \node[usecase] (PU8) at (8, -11) {PU8: Odbiór paczek\\przez kuriera};
    
    \node[usecase] (PU6) at (8, -15) {PU6: Tryb serwisowy};
    \node[usecase] (PU7) at (14, -15) {PU7: Monitorowanie\\systemu};

    % --- INCLUDES (Oś X = 14) ---
    \node[usecase, fill=yellow!10] (ObslTerm) at (14, 6) {Obsługa\\terminala};
    \node[usecase, fill=yellow!10] (AuthUser) at (14, 2) {Autoryzacja\\użytkownika};
    \node[usecase, fill=yellow!10] (GenKod) at (14, -5) {Generowanie\\kodu nadania};
    
    % Współdzielona autoryzacja dla PU3 i PU8
    \node[usecase, fill=yellow!10] (AuthKurier) at (14, -9.5) {Autoryzacja\\kuriera};

    % --- EXTENDS (Oś X = 2) ---
    \node[usecase, fill=red!10] (OdmPlat) at (2, 6) {Odmowa\\płatności};
    \node[usecase, fill=red!10] (BrakSkryt) at (2, -5) {Brak wolnych\\skrytek};
    
    % Współdzielone otwarcie wielu skrytek dla PU3 i PU8
    \node[usecase, fill=red!10] (OtwWielu) at (2, -9.5) {Otwarcie wielu\\skrytek};
    
    \node[usecase, fill=red!10] (Restart) at (2, -15) {Restart\\systemu};

\end{scope}

% Rysowanie ramki systemu
\scoped[on background layer]
  \node[draw, thick, inner sep=25pt, fit=(system), label=above:\textbf{System Paczkomatu}] {};

% ==========================================
% POWIĄZANIA AKTORÓW (Asocjacje)
% ==========================================
\draw[assoc] (User) -- (PU1);
\draw[assoc] (User) -- (PU2);
\draw[assoc] (User) -- (PU4);
\draw[assoc] (User) to[bend left=10] (PU5);

\draw[assoc] (Kurier) -- (PU3);
\draw[assoc] (Kurier) -- (PU8);

\draw[assoc] (Serwisant) -- (PU6);
% Wygięcie strzałki serwisanta tak, aby omijała węzeł PU6
\draw[assoc] (Serwisant) to[bend right=25] (PU7);

\draw[assoc] (SysPlat) -- (PU4);

\draw[assoc] (SysCentr) to[bend right=15] (PU1);
\draw[assoc] (SysCentr) to[bend right=5] (PU2);
\draw[assoc] (SysCentr) to[bend left=5] (PU3);
\draw[assoc] (SysCentr) -- (PU5);
\draw[assoc] (SysCentr) to[bend left=15] (PU8);
% Wygięcie do Monitorowania, aby nie przecinać węzłów Includes
\draw[assoc] (SysCentr) to[bend left=20] (PU7);

% ==========================================
% RELACJE INCLUDE
% ==========================================
\draw[include] (PU4) -- node[above, sloped, font=\footnotesize] {<<include>>} (ObslTerm);
\draw[include] (PU1) -- node[above, sloped, font=\footnotesize] {<<include>>} (AuthUser);
\draw[include] (PU2) -- node[above, sloped, font=\footnotesize] {<<include>>} (GenKod);
\draw[include] (PU3) -- node[above, sloped, font=\footnotesize] {<<include>>} (AuthKurier);
\draw[include] (PU8) -- node[below, sloped, font=\footnotesize] {<<include>>} (AuthKurier);
\draw[include] (PU6) -- node[above, sloped, font=\footnotesize] {<<include>>} (PU7);

% ==========================================
% RELACJE EXTEND
% ==========================================
\draw[extend] (PU4) -- node[right, font=\footnotesize] {<<extend>>} (PU1);
\draw[extend] (OdmPlat) -- node[above, sloped, font=\footnotesize] {<<extend>>} (PU4);
\draw[extend] (BrakSkryt) -- node[above, sloped, font=\footnotesize] {<<extend>>} (PU2);
\draw[extend] (OtwWielu) -- node[above, sloped, font=\footnotesize] {<<extend>>} (PU3);
\draw[extend] (OtwWielu) -- node[below, sloped, font=\footnotesize] {<<extend>>} (PU8);
\draw[extend] (Restart) -- node[above, sloped, font=\footnotesize] {<<extend>>} (PU6);

\end{tikzpicture}
}
\end{center}


%%%%%%%%%%%%%%%%%%%%%%%%%%%%%%%%%%%%%%%%%%%%%%%%%%
\subsection{PU1 – Odbiór paczki}

\begin{tabular}{|p{6cm}|p{4cm}|p{5cm}|}
\hline
\textbf{Nazwa PU:} Odbiór paczki &
\textbf{Numer PU:} PU1 &
\textbf{Priorytet:} wysoki \\
\hline
\textbf{Aktor podstawowy:} Użytkownik &
\multicolumn{2}{p{9cm}|}{\textbf{Typ opisu:} szczegółowy} \\
\hline
\multicolumn{3}{|p{15cm}|}{
\textbf{Udziałowcy i cele:}
\begin{itemize}
\item Użytkownik – chce odebrać paczkę
\item System centralny – potwierdzenie odbioru
\end{itemize}
} \\
\hline
\textbf{Wyzwalacz:} Użytkownik wprowadza kod lub skanuje QR &
\multicolumn{2}{p{9cm}|}{\textbf{Typ wyzwalacza:} zewnętrzny} \\
\hline
\multicolumn{3}{|p{15cm}|}{
\textbf{Powiązania:}
\begin{itemize}
\item Asocjacja: Nadanie paczki (PU2), Płatność przy odbiorze (PU5)
\item Zawieranie: Autoryzacja użytkownika
\item Rozszerzenie: Płatność przy odbiorze
\item Generalizacja: —
\end{itemize}
} \\
\hline
\multicolumn{3}{|p{15cm}|}{
\textbf{Zwykły przepływ zdarzeń:}
\begin{enumerate}
\item Użytkownik wybiera opcję „Odbiór paczki"
\item Wprowadza kod lub skanuje QR
\item System weryfikuje dane
\item System sprawdza, czy wymagana jest płatność
\item Skrytka zostaje otwarta
\item Użytkownik odbiera paczkę
\item System rejestruje odbiór i zamyka skrytkę
\end{enumerate}
} \\
\hline
\multicolumn{3}{|p{15cm}|}{
\textbf{Przepływy poboczne:}
\begin{itemize}
\item 1a. Użytkownik anuluje operację
\end{itemize}
} \\
\hline
\multicolumn{3}{|p{15cm}|}{
\textbf{Przepływy alternatywne/wyjątkowe:}
\begin{itemize}
\item Błędny kod → dostęp zablokowany
\item Brak połączenia z systemem → operacja niemożliwa
\end{itemize}
} \\
\hline
\end{tabular}

\vspace{1cm}

%%%%%%%%%%%%%%%%%%%%%%%%%%%%%%%%%%%%%%%%%%%%%%%%%%
\subsection{PU2 – Nadanie paczki}

\begin{tabular}{|p{6cm}|p{4cm}|p{5cm}|}
\hline
\textbf{Nazwa PU:} Nadanie paczki &
\textbf{Numer PU:} PU2 &
\textbf{Priorytet:} wysoki \\
\hline
\textbf{Aktor podstawowy:} Użytkownik &
\multicolumn{2}{p{9cm}|}{\textbf{Typ opisu:} szczegółowy} \\
\hline
\multicolumn{3}{|p{15cm}|}{
\textbf{Udziałowcy i cele:}
\begin{itemize}
\item Użytkownik – chce nadać paczkę
\item System centralny – rejestracja przesyłki
\end{itemize}
} \\
\hline
\textbf{Wyzwalacz:} Wybór opcji „Nadaj paczkę" &
\multicolumn{2}{p{9cm}|}{\textbf{Typ wyzwalacza:} zewnętrzny} \\
\hline
\multicolumn{3}{|p{15cm}|}{
\textbf{Powiązania:}
\begin{itemize}
\item Asocjacja: Odbiór przez kuriera (PU3)
\item Zawieranie: Generowanie kodu nadania
\item Rozszerzenie: Brak wolnych skrytek
\item Generalizacja: —
\end{itemize}
} \\
\hline
\multicolumn{3}{|p{15cm}|}{
\textbf{Zwykły przepływ zdarzeń:}
\begin{enumerate}
\item Użytkownik wybiera „Nadaj paczkę"
\item Wprowadza dane lub używa aplikacji
\item System przydziela skrytkę
\item Skrytka otwiera się
\item Użytkownik umieszcza paczkę
\item System zamyka skrytkę i generuje kod
\end{enumerate}
} \\
\hline
\multicolumn{3}{|p{15cm}|}{
\textbf{Przepływy poboczne:}
\begin{itemize}
\item 1a. Użytkownik rezygnuje
\end{itemize}
} \\
\hline
\multicolumn{3}{|p{15cm}|}{
\textbf{Przepływy alternatywne/wyjątkowe:}
\begin{itemize}
\item Brak wolnych skrytek → operacja przerwana
\end{itemize}
} \\
\hline
\end{tabular}

\vspace{1cm}

%%%%%%%%%%%%%%%%%%%%%%%%%%%%%%%%%%%%%%%%%%%%%%%%%%
\subsection{PU3 – Dostarczenie paczki przez kuriera}

\begin{tabular}{|p{6cm}|p{4cm}|p{5cm}|}
\hline
\textbf{Nazwa PU:} Dostarczenie paczki &
\textbf{Numer PU:} PU3 &
\textbf{Priorytet:} wysoki \\
\hline
\textbf{Aktor podstawowy:} Kurier &
\multicolumn{2}{p{9cm}|}{\textbf{Typ opisu:} szczegółowy} \\
\hline
\multicolumn{3}{|p{15cm}|}{
\textbf{Udziałowcy i cele:}
\begin{itemize}
\item Kurier – chce dostarczyć paczki
\item System centralny – aktualizacja statusów
\end{itemize}
} \\
\hline
\textbf{Wyzwalacz:} Logowanie kuriera &
\multicolumn{2}{p{9cm}|}{\textbf{Typ wyzwalacza:} zewnętrzny} \\
\hline
\multicolumn{3}{|p{15cm}|}{
\textbf{Powiązania:}
\begin{itemize}
\item Asocjacja: Odbiór paczki (PU1)
\item Zawieranie: Autoryzacja kuriera
\item Rozszerzenie: Otwarcie wielu skrytek
\item Generalizacja: —
\end{itemize}
} \\
\hline
\multicolumn{3}{|p{15cm}|}{
\textbf{Zwykły przepływ zdarzeń:}
\begin{enumerate}
\item Kurier loguje się
\item System weryfikuje dane
\item System otwiera odpowiednie skrytki
\item Kurier umieszcza paczki
\item System zamyka skrytki
\item Dane są wysyłane do systemu centralnego
\end{enumerate}
} \\
\hline
\multicolumn{3}{|p{15cm}|}{
\textbf{Przepływy poboczne:}
\begin{itemize}
\item Kurier pomija skrytkę
\end{itemize}
} \\
\hline
\multicolumn{3}{|p{15cm}|}{
\textbf{Przepływy alternatywne/wyjątkowe:}
\begin{itemize}
\item Błąd autoryzacji → brak dostępu
\end{itemize}
} \\
\hline
\end{tabular}

\vspace{1cm}

%%%%%%%%%%%%%%%%%%%%%%%%%%%%%%%%%%%%%%%%%%%%%%%%%%
\subsection{PU4 – Odbiór paczek przez kuriera}

\begin{tabular}{|p{6cm}|p{4cm}|p{5cm}|}
\hline
\textbf{Nazwa PU:} Odbiór paczek przez kuriera &
\textbf{Numer PU:} PU4 &
\textbf{Priorytet:} wysoki \\
\hline
\textbf{Aktor podstawowy:} Kurier &
\multicolumn{2}{p{9cm}|}{\textbf{Typ opisu:} szczegółowy} \\
\hline
\multicolumn{3}{|p{15cm}|}{
\textbf{Udziałowcy i cele:}
\begin{itemize}
\item Kurier – chce odebrać nadane paczki
\item System centralny – aktualizacja statusów przesyłek
\end{itemize}
} \\
\hline
\textbf{Wyzwalacz:} Logowanie kuriera i wybór opcji odbioru paczek &
\multicolumn{2}{p{9cm}|}{\textbf{Typ wyzwalacza:} zewnętrzny} \\
\hline
\multicolumn{3}{|p{15cm}|}{
\textbf{Powiązania:}
\begin{itemize}
\item Asocjacja: Nadanie paczki (PU2)
\item Zawieranie: Autoryzacja kuriera
\item Rozszerzenie: Otwarcie wielu skrytek
\item Generalizacja: —
\end{itemize}
} \\
\hline
\multicolumn{3}{|p{15cm}|}{
\textbf{Zwykły przepływ zdarzeń:}
\begin{enumerate}
\item Kurier loguje się do systemu
\item System weryfikuje dane kuriera
\item System identyfikuje skrytki z paczkami do odbioru
\item System otwiera odpowiednie skrytki
\item Kurier odbiera paczki
\item System zamyka skrytki
\item System aktualizuje status przesyłek w systemie centralnym
\end{enumerate}
} \\
\hline
\multicolumn{3}{|p{15cm}|}{
\textbf{Przepływy poboczne:}
\begin{itemize}
\item Kurier pomija wybrane skrytki
\end{itemize}
} \\
\hline
\multicolumn{3}{|p{15cm}|}{
\textbf{Przepływy alternatywne/wyjątkowe:}
\begin{itemize}
\item Brak paczek do odbioru → brak otwierania skrytek
\item Błąd autoryzacji → brak dostępu
\end{itemize}
} \\
\hline
\end{tabular}

\vspace{1cm}

%%%%%%%%%%%%%%%%%%%%%%%%%%%%%%%%%%%%%%%%%%%%%%%%%%
\subsection{PU5 – Płatność przy odbiorze}

\begin{tabular}{|p{6cm}|p{4cm}|p{5cm}|}
\hline
\textbf{Nazwa PU:} Płatność przy odbiorze &
\textbf{Numer PU:} PU5 &
\textbf{Priorytet:} średni \\
\hline
\textbf{Aktor podstawowy:} Użytkownik &
\multicolumn{2}{p{9cm}|}{\textbf{Typ opisu:} szczegółowy} \\
\hline
\multicolumn{3}{|p{15cm}|}{
\textbf{Udziałowcy i cele:}
\begin{itemize}
\item Użytkownik – chce zapłacić i odebrać paczkę
\item System płatniczy – realizacja transakcji
\end{itemize}
} \\
\hline
\textbf{Wyzwalacz:} Wykrycie konieczności płatności &
\multicolumn{2}{p{9cm}|}{\textbf{Typ wyzwalacza:} wewnętrzny} \\
\hline
\multicolumn{3}{|p{15cm}|}{
\textbf{Powiązania:}
\begin{itemize}
\item Asocjacja: Odbiór paczki (PU1)
\item Zawieranie: Obsługa terminala
\item Rozszerzenie: Odmowa płatności
\item Generalizacja: —
\end{itemize}
} \\
\hline
\multicolumn{3}{|p{15cm}|}{
\textbf{Zwykły przepływ zdarzeń:}
\begin{enumerate}
\item System informuje o konieczności płatności
\item Użytkownik dokonuje płatności
\item System potwierdza transakcję
\item Skrytka zostaje otwarta
\end{enumerate}
} \\
\hline
\multicolumn{3}{|p{15cm}|}{
\textbf{Przepływy poboczne:}
\begin{itemize}
\item Anulowanie płatności
\end{itemize}
} \\
\hline
\multicolumn{3}{|p{15cm}|}{
\textbf{Przepływy alternatywne/wyjątkowe:}
\begin{itemize}
\item Nieudana płatność → brak wydania paczki
\end{itemize}
} \\
\hline
\end{tabular}

\vspace{1cm}

%%%%%%%%%%%%%%%%%%%%%%%%%%%%%%%%%%%%%%%%%%%%%%%%%%
\subsection{PU6 – Wysyłanie powiadomień}

\begin{tabular}{|p{6cm}|p{4cm}|p{5cm}|}
\hline
\textbf{Nazwa PU:} Wysyłanie powiadomień &
\textbf{Numer PU:} PU6 &
\textbf{Priorytet:} średni \\
\hline
\textbf{Aktor podstawowy:} System centralny &
\multicolumn{2}{p{9cm}|}{\textbf{Typ opisu:} ogólny} \\
\hline
\multicolumn{3}{|p{15cm}|}{
\textbf{Udziałowcy i cele:}
\begin{itemize}
\item Użytkownik – otrzymuje informacje o statusie paczki
\item System centralny – informowanie o zdarzeniach
\end{itemize}
} \\
\hline
\textbf{Wyzwalacz:} Zmiana statusu paczki &
\multicolumn{2}{p{9cm}|}{\textbf{Typ wyzwalacza:} wewnętrzny} \\
\hline
\multicolumn{3}{|p{15cm}|}{
\textbf{Powiązania:}
\begin{itemize}
\item Asocjacja: PU1, PU2, PU3
\item Zawieranie: —
\item Rozszerzenie: —
\item Generalizacja: —
\end{itemize}
} \\
\hline
\multicolumn{3}{|p{15cm}|}{
\textbf{Zwykły przepływ zdarzeń:}
\begin{enumerate}
\item System wykrywa zmianę statusu paczki
\item Generuje powiadomienie
\item Wysyła powiadomienie do użytkownika
\end{enumerate}
} \\
\hline
\multicolumn{3}{|p{15cm}|}{
\textbf{Przepływy poboczne:}
\begin{itemize}
\item Opóźnienie wysłania powiadomienia
\end{itemize}
} \\
\hline
\multicolumn{3}{|p{15cm}|}{
\textbf{Przepływy alternatywne/wyjątkowe:}
\begin{itemize}
\item Brak połączenia z serwerem → powiadomienie niewysłane
\end{itemize}
} \\
\hline
\end{tabular}

\vspace{1cm}

%%%%%%%%%%%%%%%%%%%%%%%%%%%%%%%%%%%%%%%%%%%%%%%%%%
\subsection{PU7 – Tryb serwisowy}

\begin{tabular}{|p{6cm}|p{4cm}|p{5cm}|}
\hline
\textbf{Nazwa PU:} Tryb serwisowy &
\textbf{Numer PU:} PU7 &
\textbf{Priorytet:} średni \\
\hline
\textbf{Aktor podstawowy:} Serwisant &
\multicolumn{2}{p{9cm}|}{\textbf{Typ opisu:} ogólny} \\
\hline
\multicolumn{3}{|p{15cm}|}{
\textbf{Udziałowcy i cele:}
\begin{itemize}
\item Serwisant – diagnoza i naprawa systemu
\end{itemize}
} \\
\hline
\textbf{Wyzwalacz:} Logowanie serwisanta &
\multicolumn{2}{p{9cm}|}{\textbf{Typ wyzwalacza:} zewnętrzny} \\
\hline
\multicolumn{3}{|p{15cm}|}{
\textbf{Powiązania:}
\begin{itemize}
\item Asocjacja: PU8 (Monitorowanie systemu)
\item Zawieranie: Monitorowanie systemu
\item Rozszerzenie: Restart systemu
\item Generalizacja: —
\end{itemize}
} \\
\hline
\multicolumn{3}{|p{15cm}|}{
\textbf{Zwykły przepływ zdarzeń:}
\begin{enumerate}
\item Serwisant loguje się do systemu
\item System weryfikuje uprawnienia
\item System udostępnia funkcje serwisowe
\item Serwisant wykonuje operacje diagnostyczne lub naprawcze
\end{enumerate}
} \\
\hline
\multicolumn{3}{|p{15cm}|}{
\textbf{Przepływy poboczne:}
\begin{itemize}
\item Wyjście z trybu serwisowego bez zmian
\end{itemize}
} \\
\hline
\multicolumn{3}{|p{15cm}|}{
\textbf{Przepływy alternatywne/wyjątkowe:}
\begin{itemize}
\item Błędne dane logowania → brak dostępu
\end{itemize}
} \\
\hline
\end{tabular}

\vspace{1cm}

%%%%%%%%%%%%%%%%%%%%%%%%%%%%%%%%%%%%%%%%%%%%%%%%%%
\subsection{PU8 – Monitorowanie systemu}

\begin{tabular}{|p{6cm}|p{4cm}|p{5cm}|}
\hline
\textbf{Nazwa PU:} Monitorowanie systemu &
\textbf{Numer PU:} PU8 &
\textbf{Priorytet:} niski \\
\hline
\textbf{Aktor podstawowy:} System centralny &
\multicolumn{2}{p{9cm}|}{\textbf{Typ opisu:} ogólny} \\
\hline
\multicolumn{3}{|p{15cm}|}{
\textbf{Udziałowcy i cele:}
\begin{itemize}
\item System centralny – kontrola stanu paczkomatu
\item Serwisant – dostęp do informacji diagnostycznych
\end{itemize}
} \\
\hline
\textbf{Wyzwalacz:} Cykliczne sprawdzenie systemu &
\multicolumn{2}{p{9cm}|}{\textbf{Typ wyzwalacza:} wewnętrzny} \\
\hline
\multicolumn{3}{|p{15cm}|}{
\textbf{Powiązania:}
\begin{itemize}
\item Asocjacja: PU7 (Tryb serwisowy)
\item Zawieranie: —
\item Rozszerzenie: —
\item Generalizacja: —
\end{itemize}
} \\
\hline
\multicolumn{3}{|p{15cm}|}{
\textbf{Zwykły przepływ zdarzeń:}
\begin{enumerate}
\item System zbiera dane z czujników i modułów
\item Analizuje stan urządzenia
\item Przesyła raport do systemu centralnego
\end{enumerate}
} \\
\hline
\multicolumn{3}{|p{15cm}|}{
\textbf{Przepływy poboczne:}
\begin{itemize}
\item Częściowa niedostępność danych
\end{itemize}
} \\
\hline
\multicolumn{3}{|p{15cm}|}{
\textbf{Przepływy alternatywne/wyjątkowe:}
\begin{itemize}
\item Awaria czujnika → niepełne dane diagnostyczne
\end{itemize}
} \\
\hline
\end{tabular}

\vspace{1cm}

%%%%%%%%%%%%%%%%%%%%%%%%%%%%%%%%%%%%%%%%%%%%%%%%%%
\section{Diagramy}

\subsection{PU1 – Odbiór paczki}

\begin{figure}[H]
    \centering
    \includegraphics[width=0.55\linewidth]{PU1 - decyzyjny.png}
    \caption{Diagram decyzyjny}
\end{figure}

\begin{figure}[H]
    \centering
    \includegraphics[width=0.4\linewidth]{PU1 - przepływu.png}
    \caption{Diagram przepływu}
\end{figure}

\begin{figure}[H]
    \centering
    \includegraphics[width=0.6\linewidth]{PU1 - stanów.png}
    \caption{Diagram stanów}
\end{figure}

\begin{figure}[H]
    \centering
    \includegraphics[width=1.0\linewidth]{PU1 - sekwencji.png}
    \caption{Diagram sekwencji}
\end{figure}


\subsection{PU2 – Nadanie paczki}

\begin{figure}[H]
    \centering
    \includegraphics[width=0.55\linewidth]{PU2 - decyzyjny.png}
    \caption{Diagram decyzyjny}
\end{figure}
\begin{figure}[H]
    \centering
    \includegraphics[width=0.4\linewidth]{PU2 - przepływu.png}
    \caption{Diagram przepływu}
\end{figure}
\begin{figure}[H]
    \centering
    \includegraphics[width=0.9\linewidth]{PU2 - stanów.png}
    \caption{Diagram stanów}
\end{figure}
\begin{figure}[H]
    \centering
    \includegraphics[width=1\linewidth]{PU2 - sekwencji.png}
    \caption{Diagram sekwencji}
\end{figure}

\subsection{PU3 – Dostarczenie paczki przez kuriera}

\begin{figure}[H]
    \centering
    \includegraphics[width=0.55\linewidth]{PU3 - decyzyjny.png}
    \caption{Diagram decyzyjny}
\end{figure}
\begin{figure}[H]
    \centering
    \includegraphics[width=0.4\linewidth]{PU3 - przepływu.png}
    \caption{Diagram przepływu}
\end{figure}
\begin{figure}[H]
    \centering
    \includegraphics[width=0.9\linewidth]{PU3 - stanów.png}
    \caption{Diagram stanów}
\end{figure}
\begin{figure}[H]
    \centering
    \includegraphics[width=1\linewidth]{PU3 - sekwencji.png}
    \caption{Diagram sekwencji}
\end{figure}

\subsection{PU4 – Odbiór paczek przez kuriera}

\begin{figure}[H]
    \centering
    \includegraphics[width=0.55\linewidth]{PU4 - decyzyjny.png}
    \caption{Diagram decyzyjny}
\end{figure}
\begin{figure}[H]
    \centering
    \includegraphics[width=0.4\linewidth]{PU4 - przepływu.png}
    \caption{Diagram przepływu}
\end{figure}
\begin{figure}[H]
    \centering
    \includegraphics[width=0.9\linewidth]{PU4 - stanów.png}
    \caption{Diagram stanów}
\end{figure}
\begin{figure}[H]
    \centering
    \includegraphics[width=1\linewidth]{PU4 - sekwencji.png}
    \caption{Diagram sekwencji}
\end{figure}

\subsection{PU5 – Płatność przy odbiorze}

\begin{figure}[H]
    \centering
    \includegraphics[width=0.55\linewidth]{PU5 - decyzyjny.png}
    \caption{Diagram decyzyjny}
\end{figure}
\begin{figure}[H]
    \centering
    \includegraphics[width=0.4\linewidth]{PU5 - przepływu.png}
    \caption{Diagram przepływu}
\end{figure}
\begin{figure}[H]
    \centering
    \includegraphics[width=0.9\linewidth]{PU5 - stanów.png}
    \caption{Diagram stanów}
\end{figure}
Diagram sekwencji dla PU5 jest częscią diagramu sekwencji PU1.

%%%%%%%%%%%%%%%%%%%%%%%%%%%%%%%%%%%%%%%%%%%%%%%%%%
\section{Specyfikacja danych systemowych}

\subsection{Wprowadzenie}
Niniejszy dokument stanowi rozszerzenie specyfikacji funkcjonalnej systemu automatycznego paczkomatu. Skupia się na definicji danych przetwarzanych przez system, ich reprezentacji fizycznej oraz relacjach zachodzących między nimi w trakcie procesów nadawania i odbioru przesyłek.

\subsection{Klasyfikacja i parametry danych}
W poniższej tabeli zestawiono kluczowe dane krążące w systemie, uwzględniając ich wymiar fizyczny, jednostki oraz precyzję przetwarzania.

\begin{table}[h!]
\small
\centering
\begin{tabularx}{\textwidth}{l l l l X}
\toprule
\textbf{Nazwa danej} & \textbf{Wymiar} & \textbf{Jednostka} & \textbf{Precyzja} & \textbf{Zakres / Format} \\ \midrule
Kod odbioru / QR  & Alfanumeryczny & -- & 1:1 & 6-12 znaków / String \\
Kwota płatności  & Waluta & PLN & 0,01 PLN & 0,00 -- 9999,99 \\
Status skrytki  & Logiczny & -- & 1 bit & \{Otwarta, Zamknięta\} \\
Zajętość skrytki  & Logiczny & -- & 1 bit & \{Wolna, Zajęta\} \\
Czas przechowywania  & Czas & h, min & 1 s & Timestamp (Unix) \\
Temperatura  & Temperatura & $^\circ$C & 0,1 $^\circ$C & -40,0 -- 85,0 $^\circ$C \\
ID Przesyłki  & Identyfikator & -- & 1:1 & UUID / Kod kreskowy \\
Uprawnienia  & Poziom dostępu & -- & Dyskretna & \{User, Kurier, Serwis\} \\
\bottomrule
\end{tabularx}
\caption{Zestawienie parametrów danych w systemie paczkomatu.}
\end{table}

\subsection{Formaty wymiany i struktury danych}
System wykorzystuje zróżnicowane formaty danych w zależności od warstwy komunikacyjnej:

\begin{itemize}
    \item \textbf{Warstwa Centralna (Cloud):} Komunikacja z systemem centralnym odbywa się za pomocą formatu \textbf{JSON}. Przesyłane są w nim informacje o statusie paczek oraz raporty monitoringu.
    \item \textbf{Warstwa Sprzętowa (Hardware):} Sygnały z czujników skrytek (wejścia) oraz sterowanie ryglem (wyjścia) to dane binarne lub ramki protokołów komunikacji szeregowej.
    \item \textbf{Struktury wewnętrzne (C/C++):}
    \begin{itemize}
        \item \texttt{struct Locker}: przechowuje ID skrytki, jej rozmiar, stan rygla oraz obecność paczki.
        \item \texttt{struct UserSession}: przechowuje dane o aktualnie zalogowanym aktorze (Użytkownik, Kurier lub Serwisant).
    \end{itemize}
\end{itemize}

\subsection{Model relacyjny danych}
W oparciu o analizę przypadków użycia (PU1–PU8), zidentyfikowano następujące relacje między encjami:

\begin{enumerate}
    \item \textbf{Paczka $\leftrightarrow$ Skrytka:} Relacja 1:1 w danym momencie czasu. Skrytka zostaje przypisana do paczki w momencie nadania (PU2) i zwolniona po odbiorze (PU1).
    \item \textbf{Kurier $\leftrightarrow$ Skrytki:} Relacja 1:N. Podczas dostarczania (PU3) kurier może otworzyć wiele skrytek jednocześnie.
    \item \textbf{Paczka $\leftrightarrow$ Płatność:} Relacja 1:1 dla przesyłek pobraniowych. Status płatności (PU5) jest warunkiem koniecznym do wyzwolenia sygnału otwarcia rygla.
    \item \textbf{System Centralny $\leftrightarrow$ Zdarzenia:} Relacja 1:N. Każda operacja (PU1, PU2, PU3) generuje rekord w logu przesyłany do systemu centralnego w celu aktualizacji statusu.
\end{enumerate}

\subsection{Wymagania techniczne i bezpieczeństwo}
Dane wrażliwe, takie jak kody odbioru oraz dane płatnicze, wymagają szyfrowania na poziomie software'u. Dane o stanie systemu (monitoring) są przesyłane cyklicznie w celu wczesnego wykrywania awarii czujników (PU8).

\subsection{Diagram komponentów}

\begin{figure}[H]
    \centering
    \includegraphics[width=0.9\linewidth]{diagram komponentów.png}
    \caption{Diagram stanów}
\end{figure}

\end{document}